\section{Slice design and parameter rationale}\label{sec:slices}
Each service class maps to one S-NSSAI. We use the standardised SST values
(1\,=\,eMBB, 2\,=\,URLLC, 3\,=\,MIoT~\cite{ts23501}). We add an explicit SD so that further
tenants can later be added under the same SST. Each S-NSSAI also gets its own DNN and IPv4
pool, so a packet's source address alone identifies its slice. That makes per-slice policing,
capture filtering and KPI attribution trivial. Table~\ref{tab:slices} lists every parameter;
the reasons for each choice follow.

\begin{table}[H]\centering\small
\caption{Slice configuration (\code{configs/slices.yaml}, \code{configs/open5gs/*.yaml},
\code{scripts/30\_subscribers.sh}).}\label{tab:slices}
\begin{tabularx}{\linewidth}{l>{\raggedright}X>{\raggedright}X>{\raggedright\arraybackslash}X}
\toprule
Parameter & \slice{eMBB} & \slice{URLLC} & \slice{mIoT}\\
\midrule
S-NSSAI (SST/SD) & 1 / 000001 & 2 / 000002 & 3 / 000003\\
DNN / UE pool & \code{internet} / 10.45.0.0/16 & \code{urllc} / 10.46.0.0/16 & \code{iot} / 10.47.0.0/16\\
User plane & UPF-A (shared), \code{ogstun} & \textbf{UPF-B (dedicated)}, \code{ogstun2} & UPF-A (shared), \code{ogstun3}\\
N3 endpoint & 10.10.0.1 & 10.10.0.3 & 10.10.0.1\\
5QI / ARP priority & 9 / 8 & 80 / 1 & 9 / 12\\
Session-AMBR DL/UL & 200 / 50 Mbit/s & 20 / 20 Mbit/s & 2 / 2 Mbit/s\\
Slice MBR (tc HTB) & 150 Mbit/s & 20 Mbit/s & 5 Mbit/s\\
Added delay / loss & 0\,ms / 0\,\% (live-tunable) & 0\,ms / 0\,\% (live-tunable) & 0\,ms / 0\,\% (live-tunable)\\
Subscribed UEs & embb1, embb2, multi & urllc1, multi & miot1, miot2\\
\bottomrule
\end{tabularx}
\end{table}

\paragraph{Dedicated vs.\ shared user plane.} The URLLC slice gets its own UPF process (UPF-B)
with its own PFCP association, GTP-U endpoint and tun device. The aim is to isolate its
forwarding path from the bursty eMBB load, both in CPU time and in queues. In a real
deployment the same choice lets URLLC be placed at the factory edge (local breakout) while
eMBB/mIoT can be centralised. mIoT carries only kbit/s, so a dedicated UPF would cost resources
for no benefit; it therefore shares UPF-A with eMBB. Experiment E5 (\S\ref{sec:e5}) tests this
decision directly by moving URLLC onto UPF-A under eMBB overload.

\paragraph{QoS parameters.} 5QI~9 is the default non-GBR class for buffered video and TCP
traffic, which fits eMBB and telemetry. URLLC uses 5QI~80. That is the standardised
\emph{non-GBR low-latency} class (10\,ms packet delay budget), and it needs no GBR flow
template. Delay-critical GBR 5QIs (82--85) would require PCC rules with GBR values, which
UERANSIM's default QoS-flow handling does not exercise. The ARP gives URLLC the highest
priority (1) and mIoT the lowest (12). Session-AMBRs are provisioned in the UDR and delivered
to the UE in the PDU Session Establishment Accept (verified with \code{nr-cli ps-list},
\S\ref{sec:e1}).

\paragraph{Why an extra tc layer.} Open5GS's UPF polices the Session-AMBR. We measured
230/57\,Mbit/s for a 200/50 eMBB AMBR and 1.9\,Mbit/s for the 2\,Mbit/s mIoT AMBR with the
tc caps lifted (\code{results/ambr\_check.txt}). However, the AMBR is \emph{per session}, is
fixed when the session is established, and cannot express delay or loss. A software RAN also
has no scheduler that could give one slice priority over another. We therefore implement a
\emph{per-slice} MBR and the delay/loss ``transport'' parameters with Linux \code{tc}: an HTB
class plus a netem child per slice, in both directions. The downlink is shaped on the slice's
UPF tun device and the uplink on the N6 veth, with one class per UE pool. Every parameter can
be changed live (\code{tc class/qdisc change}) without touching any UE session. This is the
``live slice parameter'' exposed by the sandbox. The binding limit is the lower of the two:
the tc MBR for eMBB (150\,Mbit/s) and URLLC (20\,Mbit/s), and the Session-AMBR for mIoT
(2\,Mbit/s). The values cover the scenario's peak demand (inspection video
$\approx$\,2$\times$50\,Mbit/s plus headroom; URLLC $\approx$100$\times$ its nominal
0.05\,Mbit/s; mIoT far above its few kbit/s), while staying low enough to protect the other
slices.

\paragraph{Slice selection path.} The gNB advertises all three S-NSSAIs in NG~Setup. The AMF
supports them in its PLMN list. The NSSF maps each S-NSSAI to an NSI served by the NRF. The
SMF registers \code{info.s\_nssai} for all three. Each subscriber's UDR record lists only its
permitted S-NSSAIs, so a UE that requests an unsubscribed slice is rejected
(\S\ref{sec:e1}).
